\makereporttitle{2025}{mai}

\makereportpreamble

\vspace{0.5em}

Luna aceasta am continuat să citesc/studiez materialele relevante temei mele de doctorat și am reușit să fac progres pe mai multe direcții pe care le abordez în cadrul doctoratului.

Între 14 și 17 mai, am fost la Iași pentru a participa la a 14-a ediție a ``\href{https://fmi.univ-ovidius.ro/en/cercetare/wyrm2025/}{\textit{Workshop for Young Researchers in Mathematics}}''. Am urmărit prezentările altor cercetători tineri, de interes atât pe zona de fizică computațională cât și de fizică matematică, și am format legături profesionale cu ei.

Între 18 și 25 mai am fost plecat în Italia, în L'Aquila, unde am participat la workshop-ul ``\href{https://trimester2025.math.gssi.it/workshop4/}{\textit{Hidden structures in dynamical systems, optimization, and machine learning}}'', organizat de către \href{https://gssi.it/}{\textit{Gran Sasso Science Institute}}. În cadrul acestuia au avut loc prelegeri despre subiecte clasice din analiză numerică, tehnici de optimizare, teoria controlului și analiză stocastică, dar și din unele mai recente, cum ar fi învățarea automată (în particular, \emph{Physics-Inspired Neural Networks}). Cu ajutorul acestei experiențe am reușit să înțeleg mai bine starea tehnicii în domeniul meu și în domenii conexe și am creat punți de comunicare, atât cu alți doctoranzi, cât și cu cercetători consacrați (din Italia, Germania, Marea Britanie, SUA).

De asemenea, la sfârșitul lunii mă voi deplasa la Timișoara pentru a participa la conferința ``\href{https://timconference.uvt.ro/}{TIM 2025}'', unde voi ține o prezentare intitulată ``\textit{Critical behavior at the transition to chaos in Hamiltonian systems}''.

În acest timp, am continuat să lucrez pe temele și proiectele mele de cercetare. Am optimizat și finalizat codul de calcul numeric pentru analiza Hamiltonianului cuadrupol bosonic cuantic, am reușit să implementez un cod performant pentru calcularea descompunerii Cholesky (ca parte a cercetării mele pe eficientizarea metodelor de rezolvare a sistemelor liniare mari) și am parcurs foarte multe materiale despre \emph{Conformal Field Theories} (CFT), un subiect din zona de teorii de câmp cuantic, pe care am început să-l explorez în luna precedentă.

\textbf{Observațiile conducătorului științific}:

\vspace{3em}

\makereportpostamble{28 mai 2025}
